\documentclass[12pt]{article}

\usepackage[utf8]{inputenc}
\usepackage[romanian]{babel}
\usepackage{amsmath, amssymb, amsthm}
\usepackage{geometry}
\geometry{a4paper, margin=2.5cm}

\title{Definiții ale Intuiției}
\author{}
\date{}

\theoremstyle{definition}
\newtheorem{definition}{Definiție}

\begin{document}

\maketitle

\begin{definition}[Intuiția în sens lexical]
În accepțiunea dicționarelor explicative, \emph{intuiția} desemnează o formă de cunoaștere imediată, spontană, care nu se bazează pe raționamente explicite. Este percepută ca o înțelegere directă a unui adevăr, a unei situații sau a unei relații.
\end{definition}

\begin{definition}[Intuiția în filozofie: Kant]
În filosofia lui Immanuel Kant, \emph{intuiția} reprezintă modul prin care mintea primește în mod nemijlocit datele sensibilității. Spațiul și timpul sunt forme ale intuiției sensibile, prin care obiectele devin posibile ca fenomene.
\end{definition}

\begin{definition}[Intuiția în filozofie: Bergson]
Pentru Henri Bergson, \emph{intuiția} este o formă superioară de cunoaștere, prin care spiritul pătrunde direct în esența lucrurilor, depășind analiza discursivă. Ea oferă acces la durata reală și la dinamica internă a realității.
\end{definition}

\begin{definition}[Intuiția în pedagogie]
În pedagogie, \emph{intuiția} este capacitatea elevului de a sesiza rapid structuri, relații sau idei matematice, înaintea unei justificări formale. Ea funcționează ca suport pentru înțelegere, anticipare, modelare și orientare în rezolvarea problemelor.
\end{definition}

\begin{definition}[Intuiția în didactica matematicii]
În didactica matematicii, \emph{intuiția} este un instrument cognitiv care facilitează trecerea de la experiența concretă la formalizare. Ea permite:
\begin{enumerate}
    \item identificarea rapidă a unei strategii de rezolvare;
    \item construirea unor modele mentale relevante;
    \item anticiparea rezultatelor;
    \item înțelegerea relațiilor matematice înaintea demonstrației;
    \item conectarea reprezentărilor vizuale, verbale și simbolice.
\end{enumerate}
\end{definition}

\end{document}
